\subsection{K3 components and their actual cyclic quotient models}

\begin{statement}{C.2.UN.24}{Proved}{K3 components are elliptic in positive algebraic dimension; algebraic dimension zero has an exact occurrence criterion}
Let $f:U\to\Delta$ be a proper holomorphic map from a
smooth complex threefold, with connected two-torus fibres
over $\Delta^*$. Let $S_i$ be an irreducible component of
its central support, and suppose that a resolution of
$S_i^\nu$ has a K3 minimal model $M$.
If $a(M)>0$, there is an actual elliptic fibration
$M\to\mathbb P^1$.

Among K3 surfaces of algebraic dimension zero, the minimal
birational classes which occur in this way are exactly the
minimal resolutions of cyclic quotients of complex two-tori
by nontrivial symplectic automorphisms without eigenvalue
one on the tangent space. After translating the origin, the generator
is linear of order $n=2,3,4$, or $6$. The quotient singularities
and exceptional configurations are exactly
\[
\begin{array}{c|c|c|c}
n&\text{quotient singularities}&
 \text{number of exceptional curves}&e(A/\langle g\rangle)\\ \hline
2&16A_1&16&8\\
3&9A_2&18&6\\
4&4A_3+6A_1&18&6\\
6&A_5+4A_2+5A_1&18&6.
\end{array}                                                    \tag{C.K3.1}
\]
Here the cyclic quotient in the criterion is required to
have algebraic dimension zero; the table applies to the
cyclic K3 quotient models in every algebraic dimension.

Every cyclic quotient K3 minimal class just described is
realized by a compact smooth threefold over $\mathbb P^1$
with connected torus general fibre and with both special
reduced supports of Euler characteristic zero. The central
support has a component whose normalization resolution has
exactly that K3 minimal model. The total space and construction
are projective when the given torus is projective.
\end{statement}

\begin{proof}
Use the actual normalized root base change and equivariant
resolution of C.2.UN.16. For the chosen component, let $E$
be a smooth strict covering component and $G$ its cyclic
stabilizer. The comparison is
\[
 E\longrightarrow S_i^\nu,
 \qquad E/G\dashrightarrow S_i^\nu\dashrightarrow M,
                                                               \tag{C.K3.2}
\]
where the first map is generically finite, the quotient
comparison is proper birational, and the remaining arrows
are resolved by their proper graphs. In particular, a
nonzero holomorphic two-form on $M$ pulls back, by its
actual Jacobian map, to a nonzero holomorphic two-form on
a resolution of the graph in (C.K3.2). Point-blowup
invariance then gives
\[
                  p_g(E)\ge1,\qquad \kappa(E)\ge0.
                                                               \tag{C.K3.3}
\]
Also $a(E)=a(M)$. To check this equality through the
specified map, a meromorphic function on $E$ has elementary
symmetric functions along the finite generic sheets.
They descend meromorphically to $M$, so that the original
function satisfies a monic polynomial over $\mathcal M(M)$.
Thus the two meromorphic function fields have the same
transcendence degree.

The two-dimensional space of fields on $E$ in C.2.UN.16
has generic rank one or two. In rank one, C.2.UN.20 gives
an actual rational or elliptic fibration on a birational
model of $M$. The rational alternative would make all
positive pluricanonical sections vanish, by restricting
them to its smooth general fibre. The pulled-back K3
form excludes it. Hence the quotient fibration is elliptic.

In rank two, C.2.UN.22 gives $\chi(\mathcal O_E)=0$ and
its precise torus or elliptic-fibration alternatives.
First suppose $a(M)=a(E)=1$. In the elliptic alternative,
the connected-fibre map $q:E'\to B$ has
\[
                     \mathcal M(E')=q^*\mathcal M(B).
                                                               \tag{C.K3.4}
\]
Indeed every meromorphic function is algebraic over any
one nonconstant function on $B$ and hence constant on
the connected general fibres; the proper graph descends
it, as in (C.RF.9). The equality makes this field
intrinsic and therefore $G$-invariant. Resolve the graph
of $q$ equivariantly and take its Stein factor. The group
acts on that factor, and the actual finite quotient is a
map to $B/G$ whose general fibre is $F/K$, with
$K=\ker(G\to\operatorname{Aut}B)$. This is the same
quotient and Stein comparison proved in C.2.UN.20.
Its general fibre is connected and rational or elliptic;
the K3 form again excludes the rational case. In the
torus alternative the minimal model of $E$ is a torus
and the cyclic quotient argument below applies.

Now suppose $a(M)=a(E)=2$. We work on the smooth
projective minimal model $E_0$, to which the finite
action descends. In rank two, (C.K3.3) and
$\chi(\mathcal O_{E_0})=0$, together with the compact
surface classification already used in C.3.1
\cite[Theorem~3.1]{Peters1990}, leave the following
possibilities. If $\kappa(E_0)=0$, it is an Abelian
surface: K3 and Enriques surfaces have coherent Euler
characteristic $2$ and $1$, and bielliptic surfaces have
$p_g=0$. If $\kappa(E_0)=1$, its elliptic Iitaka
fibration is intrinsic. More explicitly, $G$ acts on
each space of pluricanonical sections, so a sufficiently
large pluricanonical map, its equivariant graph resolution,
and its Stein factor are all $G$-equivariant. The
connected general fibres are elliptic. Taking the
quotient gives the actual fibration with general fibre
$F/K$ just described. It is elliptic because its total
quotient has the K3 minimal model. Thus the only
remaining positive-dimensional case is a cyclic quotient
of a two-torus.

For use in all these branches, an elliptic fibration on
a smooth birational model of $M$ descends to $M$ itself.
Take an actual common resolution
\[
                  M\xleftarrow{\rho}X\xrightarrow{q}B
                                                               \tag{C.K3.5}
\]
with connected general elliptic fibre $F$. Since $K_M=0$,
the point-blowup formula gives
$K_X=\sum_j a_jE_j$ with $a_j>0$ on every
$\rho$-exceptional prime curve. Adjunction gives
$K_X\cdot F=0$. Each $E_j\cdot F\ge0$, so all these
intersections are zero and every exceptional curve is
vertical for $q$. Every connected fibre of $\rho$ is
therefore sent to a point. The proper graph, together
with $\rho_*\mathcal O_X=\mathcal O_M$, descends $q$
to a holomorphic map $M\to B$ with connected fibres.
The Leray injection
$H^1(B,\mathcal O_B)\hookrightarrow H^1(M,\mathcal O_M)=0$
gives $B\simeq\mathbb P^1$.

We next compute the cyclic quotient branch. By
C.2.UN.23, a torus minimal model $A$ of $E$ has an
actual holomorphic cyclic action whose quotient is
birational to $M$. The K3 alternative there says that
its generator $g(z)=Lz+b$ has no eigenvalue one and
$\det L=1$. Since $I-L:A\to A$ is an isogeny, the
equation $(I-L)p=b$ has a solution. Translation by
$p$ conjugates $g$ to the linear automorphism $L$.
Write its eigenvalues as $\lambda,\lambda^{-1}$,
with $\lambda$ a primitive $n$th root. The integral
action on $H_1(A,\mathbb Z)$ has characteristic polynomial
\[
 (x^2-sx+1)^2
   =x^4-2sx^3+(s^2+2)x^2-2sx+1,
 \qquad s=\lambda+\lambda^{-1}\in[-2,2].
                                                               \tag{C.K3.6}
\]
Its integral coefficients imply $2s\in\mathbb Z$
and $s^2\in\mathbb Z$. Writing $s=k/2$ shows that
$k$ is even, hence $s\in\mathbb Z$. The value $s=2$
is excluded by the absence of eigenvalue one.
Thus $s=-2,-1,0,1$, giving precisely
$n=2,3,4,6$.

For every nonidentity power, its fixed points are the
kernel of an isogeny and their exact number is
\[
 \#\operatorname{Fix}(g^k)
 =\det_{\mathbb R}(I-L^k)
 =|\det_{\mathbb C}(I-L^k)|^2
 =\bigl(2-2\cos(2\pi k/n)\bigr)^2.
                                                               \tag{C.K3.7}
\]
At a point with stabilizer of order $d$, the local
action is $(u,v)\mapsto(\zeta_du,\zeta_d^{-1}v)$.
Its invariant coordinates $x=u^d,y=v^d,z=uv$ give
exactly $xy=z^d$. The height-one toric resolution
in (C.CQ.11) gives its chain of $d-1$ rational
$(-2)$-curves. Formula (C.K3.7) now gives the table
by subtracting points with larger stabilizers and
dividing by the remaining orbit sizes. For $n=4$
there are $4$ full-stabilizer points and
$(16-4)/2=6$ order-two orbits. For $n=6$ the
counts are $1$, $(9-1)/2=4$, and
$(16-1)/3=5$ for stabilizer orders $6,3,2$.
The order-two and order-three cases have $16$ and
$9$ fixed points respectively. The finite quotient
Euler formula gives the last column; adding the
exceptional ranks gives $24$ in every row.
The computation is reproduced with exact integral
matrices in \texttt{code/D\_cyclic\_k3.py}.

If this quotient K3 $M$ is projective, the table gives
an explicit elliptic pencil criterion. Choose one
exceptional $(-2)$-curve from each of four different
singular points, and call them $C_1,\ldots,C_4$.
They are pairwise disjoint. For an ample integral
class $H$ on $M$, put
\[
 H'=H+\sum_{i=1}^4\frac{H\cdot C_i}{2}C_i,
 \quad H'\cdot C_i=0,
 \quad H'^2=H^2+\frac12\sum_{i=1}^4(H\cdot C_i)^2>0.
                                                               \tag{C.K3.8}
\]
Lagrange's four-square theorem
\cite[Theorem~II, Corollary~II]{Lagrange1770}
expresses the positive rational number $H'^2/2$ as
$\sum_i x_i^2$ with $x_i\in\mathbb Q$: for a
positive fraction $a/b$, apply the integral theorem
to $ab$ and divide the four summands by $b$.
Consequently
\[
 D=H'-\sum_{i=1}^4x_iC_i,
 \qquad D^2=0,
 \qquad D\cdot H'=H'^2>0.
                                                               \tag{C.K3.9}
\]
Clear denominators to obtain a nonzero isotropic
integral divisor class. The original elliptic-pencil
criterion \cite[\S3, Corollary~3, with the reflection
argument of \S6]{PiatetskiShapiroShafarevich1971}
then produces an actual base-point-free elliptic
pencil on $M$. This proves the remaining projective
branch. If the quotient has algebraic dimension
one, then $a(A)=1$ as well. Choose a nonconstant
meromorphic function on $A$, and let $D$ be a
general member of its pencil after removing fixed
components. Its self-intersection is nonnegative,
by intersecting two distinct members. In fact
$D^2=0$. To verify this, average the positive
integration current $[D]$ under translations of
the compact torus. The result is the constant
semipositive $(1,1)$-form representing $c_1(D)$.
If $D^2>0$, this form is positive definite, so
the integral class is a polarization and the
torus is projective, contrary to $a(A)=1$.
The equality $D^2=0$ shows that two distinct
members have no common point. The pencil thus
defines a holomorphic map, and its Stein factor
$q:A\to B$ has connected smooth general fibre
with trivial normal line. Adjunction with
$K_A=0$ gives genus one. The field identity
(C.K3.4) makes this fibration $G$-equivariant;
its quotient and (C.K3.5) give the actual
elliptic pencil on $M$. This completes the
assertion for $a(M)>0$.

If $a(M)=0$, equality of algebraic dimensions in
(C.K3.2) gives $a(E)=0$. The final assertion of
C.2.UN.16 identifies its minimal model as a torus
or a Hopf surface. The nonzero canonical form
in (C.K3.3) excludes Hopf. Thus (C.K3.2) and
C.2.UN.23 give exactly the cyclic symplectic torus
quotient of (C.K3.6)--(C.K3.7). This proves necessity
of the stated algebraic-dimension-zero criterion.

Finally we prove the occurrence assertion with its
specified compact comparison. Start with any torus
$A$ and one of the actual cyclic symplectic actions
above, of order $n$. Form the finite analytic quotient
\[
 W=(A\times\mathbb P^1_s)/
       \langle(a,s)\mapsto(g(a),\zeta_ns)\rangle,
 \qquad W\longrightarrow\mathbb P^1_t,
 \qquad t=s^n.
                                                               \tag{C.K3.10}
\]
Over $t\ne0,\infty$, the $n$ sheets of the base
change are permuted transitively, and their quotient
fibre is exactly $A$. The reduced fibres over zero
and infinity are exactly $A/\langle g\rangle$.
Choose an analytic resolution of the total quotient
and its two central supports, isomorphic over the
punctured base and at a generic point of each of
these two original components. The strict transform
of the latter maps properly birationally to the
specified torus quotient. Its normalization resolution
therefore has the K3 minimal model already proved
in C.2.UN.23. The resulting smooth compact total
space has connected fibres: its Stein factor over
$\mathbb P^1$ has generic degree one and hence is
the identity finite map to the normal base.

Each of its two reduced special supports has even
Euler characteristic by C.2.UN.1. Apply the actual
vertical point and quartic-curve blowup operations
of C.2.UN.4 separately over zero and infinity.
They change the reduced Euler number by $+2$ or
$-2$ until it is zero, preserve the punctured torus
family, and leave the original strict component
birational to the same torus quotient. Resolution
and these blowups are projective morphisms when
$A$ is projective. This supplies all asserted
compact realizations, including every class of
algebraic dimension zero in the criterion.

The algebraic-dimension-zero occurrence criterion
is thus an equivalence. In positive algebraic
dimension the theorem gives the elliptic-fibration
restriction and realizes its cyclic quotient
subfamily by the explicit compact construction.
\end{proof}
